\section{Сравнительный анализ нейросетевых архитектур и проектирование методики обнаружения}
\label{ch:design}

\noindent
\textit{Содержание главы соответствует второму этапу постановки
задачи (см.~ВКР, Этап~2): анализ и сравнение архитектур нейронных
сетей; формирование теоретической модели и требований к данным;
генерация и обоснование синтетических тестовых сценариев.}\\

\subsection{Принципы сравнения и методологические замечания}

Цель главы --- обоснование выбора нейросетевой архитектуры и
формализация методики обработки данных и тестирования.
Сравнение архитектур проводится в два этапа: \emph{(а)} теоретический
анализ опубликованных в литературе NIDS-результатов с явной привязкой
к первоисточникам (см. \tabref{tab:lit-results}); \emph{(б)}
собственный эксперимент в равных условиях, описываемый на
этапе~\ref{ch:experiments}.

\textbf{Методологическое замечание.} Прямое количественное сравнение
результатов из разных литературных источников некорректно из-за:
(а)~различий датасетов (NSL-KDD, KDDCup99, UNSW-NB15, CICIDS2017);
(б)~различий схем препроцессинга и разделения выборок;
(в)~различной интерпретации метрик при дисбалансе
классов~\cite{khraisat2019survey, sommer2010outside, pendlebury2019tesseract}.
Поэтому литературные результаты используются как качественные
индикаторы.

\subsection{Опубликованные результаты применения архитектур в NIDS}

В \tabref{tab:lit-results} сведены опубликованные в рецензируемых
источниках NIDS-результаты с указанием датасета и постановки.

{
\renewcommand{\arraystretch}{1.2}
\setlength{\tabcolsep}{3pt}
\small
\begin{longtable}{@{}p{3.0cm}p{2.4cm}p{2.6cm}p{1.8cm}p{4.7cm}@{}}
\caption{Опубликованные результаты применения нейронных сетей в задачах
NIDS}\label{tab:lit-results}\\
\toprule
\textbf{Источник} & \textbf{Архитектура} & \textbf{Датасет} &
\textbf{Задача} & \textbf{Результат} \\
\midrule
\endfirsthead
\multicolumn{5}{@{}l@{}}{\textit{Продолжение таблицы~\ref{tab:lit-results}}}\\
\toprule
\textbf{Источник} & \textbf{Архитектура} & \textbf{Датасет} &
\textbf{Задача} & \textbf{Результат} \\
\midrule
\endhead
\bottomrule
\endfoot
Kim и соавт.~\cite{kim2016lstm} & LSTM-RNN & KDD Cup 99 & binary &
Accuracy~$\approx 96{,}93\%$; DR~$\approx 98{,}88\%$;
FAR~$\approx 10{,}04\%$ \\
\addlinespace
Yin и соавт.~\cite{yin2017dlrnn} & RNN & NSL-KDD (KDDTest+) & binary &
Accuracy~$\approx 83{,}28\%$ при 80 hidden units \\
\addlinespace
Yin и соавт.~\cite{yin2017dlrnn} & RNN & NSL-KDD (KDDTest+) & multiclass &
Accuracy~$\approx 81{,}29\%$ \\
\addlinespace
Moustafa и Slay~\cite{moustafa2016eval} & ANN & UNSW-NB15 & binary &
Accuracy~$\approx 81{,}3\%$; FAR~$\approx 18{,}4\%$ \\
\addlinespace
Vinayakumar и соавт.~\cite{vinayakumar2019dl} &
DNN, CNN, LSTM, CNN-LSTM &
NSL-KDD, KDDCup99, UNSW-NB15, CICIDS2017 &
binary, multiclass &
CNN-LSTM показал наибольшие F1/Accuracy среди DL-семейства данной
работы на UNSW-NB15 и CICIDS2017 \\
\addlinespace
Mirsky и соавт.~\cite{mirsky2018kitsune} &
Kitsune (ансамбль AE) &
собственный стенд (Mirai, SSDP Flood и др.) &
semi-supervised, online &
AUC: Mirai~$\approx 0{,}9988$; SSDP~$\approx 0{,}9844$;
ARP MitM~$\approx 0{,}9606$ \\
\addlinespace
Sharafaldin и соавт.~\cite{sharafaldin2018toward} & MLP, RF, KNN &
CICIDS2017 & multiclass &
MLP: Precision~$\approx 0{,}77$, Recall~$\approx 0{,}83$;
RF/KNN превосходят MLP в той же работе \\
\addlinespace
Ahmad и соавт.~\cite{ahmad2021nids} (обзор) &
LSTM, BiLSTM, CNN-LSTM & NSL-KDD, UNSW-NB15, CICIDS2017 &
binary, multiclass &
BiLSTM и CNN-LSTM устойчиво превосходят LSTM по Accuracy/F1 \\
\end{longtable}
}

Выявленные тенденции в литературе:

\begin{enumerate}
    \item Гибридные модели (CNN-LSTM, BiLSTM) устойчиво превосходят
    как чисто свёрточные, так и чисто рекуррентные модели по F1 в
    \emph{внутренних} сравнениях соответствующих работ.
    \item Высокие показатели на KDD-семействе (\(>95\,\%\) Accuracy)
    не должны интерпретироваться как доказательство практической
    пригодности из-за известных дефектов
    датасета~\cite{tavallaee2009nslkdd, sommer2010outside}.
    \item На современных датасетах (UNSW-NB15, CICIDS2017) разрыв
    между классическими ML- и DL-моделями существенно
    меньше~\cite{ferrag2020dlcyber, ahmad2021nids}.
\end{enumerate}

\subsection{Обоснование выбора CNN-LSTM как основной архитектуры}
\label{subsec:cnn-lstm-rationale}

С учётом теоретического анализа в качестве основной
(\emph{proposed}) архитектуры выбирается \textbf{CNN-LSTM}. Решение
обосновано следующими аргументами.

\begin{enumerate}
    \item \textbf{Соответствие природе данных.} Атаки в
    flow-представлении проявляются одновременно как локальные
    паттерны (например, специфическое сочетание признаков
    SYN-флуда в одном flow) и как временные тренды
    (последовательность подобных flow). CNN-LSTM моделирует обе
    компоненты в одной архитектуре.
    \item \textbf{Опубликованные результаты.} В нескольких
    независимых рецензируемых работах~\cite{vinayakumar2019dl,
    ferrag2020dlcyber, liu2019survey, ahmad2021nids} CNN-LSTM
    демонстрирует наибольшие F1/Accuracy в внутренних сравнениях
    на UNSW-NB15 и CICIDS2017.
    \item \textbf{Воспроизводимость.} Архитектура хорошо
    документирована, реализуется стандартными примитивами PyTorch,
    не требует специализированных техник обучения.
    \item \textbf{Эксплуатационная стоимость.} Inference latency
    укладывается в диапазон, приемлемый для NTA-pipeline.
    \item \textbf{Альтернативы и контрольные модели.} В качестве
    альтернатив в собственном эксперименте сохраняются BiLSTM
    (двунаправленный контекст) и Transformer-encoder (attention);
    в качестве контрольных semi-supervised baselines --- AE и VAE;
    в качестве классических baseline --- Logistic Regression,
    Random Forest, XGBoost, Isolation Forest, One-Class SVM.
\end{enumerate}

\subsection{Теоретическая модель данных и pipeline обработки}

\subsubsection{Схема признаков UNSW-NB15}

UNSW-NB15 (\emph{cleaned} распределение) содержит 44 содержательных
признака, разделяемых на пять групп
(\tabref{tab:unsw-feature-groups}).

\begin{table}[H]
\centering
\caption{Группы признаков UNSW-NB15 и их назначение}
\label{tab:unsw-feature-groups}
\renewcommand{\arraystretch}{1.2}
\begin{tabularx}{\linewidth}{@{}lXX@{}}
\toprule
\textbf{Группа} & \textbf{Состав} & \textbf{Роль в детекции} \\
\midrule
Базовые потоковые &
\emph{dur, proto, service, state, sbytes, dbytes, sttl, dttl,
sloss, dloss, rate}
&
длительность, протокол, состояние сессии, байты и пакеты в обе
стороны, общая скорость потока \\
\addlinespace
Контентные &
\emph{sload, dload, swin, dwin, smean, dmean, stcpb, dtcpb,
trans\_depth, response\_body\_len, spkts, dpkts} &
скорости, окна TCP, базовые транзакционные характеристики \\
\addlinespace
Временные &
\emph{sjit, djit, sinpkt, dinpkt, tcprtt, synack, ackdat} &
межпакетные интервалы, джиттеры, RTT, timing-сигнатуры \\
\addlinespace
Дополнительные &
\emph{ct\_state\_ttl, ct\_flw\_http\_mthd, is\_ftp\_login,
ct\_ftp\_cmd, is\_sm\_ips\_ports} &
агрегаты по протокольному поведению, бинарные флаги \\
\addlinespace
Контекстные \emph{ct\_*}&
\emph{ct\_srv\_src, ct\_srv\_dst, ct\_dst\_ltm, ct\_src\_ltm,
ct\_src\_dport\_ltm, ct\_dst\_sport\_ltm, ct\_dst\_src\_ltm} &
агрегаты по последним 100 соединениям; отражают режим
сканирования и фан-аут \\
\bottomrule
\end{tabularx}
\end{table}

\subsubsection{Pipeline обработки}

Полный pipeline принимает на вход поток flow-записей:
\[
\boxed{
\begin{array}{l}
\text{IPFIX/CSV}\ \to\ \text{Schema}\ \to\ \text{Cleanup}\ \to\ \text{Encoding}\\[2pt]
\to\ \text{Scaling}\ \to\ \text{Windowing}\ \to\ \text{Model}\\[2pt]
\to\ \text{Threshold}\ \to\ \text{Alert}
\end{array}
}
\]

\textbf{Очистка.} Применяются: замена нечисловых пропусков константой,
числовых --- нулём; унификация регистра и пробелов в
\emph{service, state, proto}; удаление дубликатов; обрезка крайних
\(0{,}1\%\) хвостов распределений. Вся статистика собирается
\emph{только} на обучающей выборке.

\textbf{Кодирование категориальных признаков.} One-hot для признаков
с малой кардинальностью (\emph{proto, state}); frequency-encoding
для \emph{service} (\(\sim 13\) уровней):
\begin{equation}
e_{\text{freq}}(c) = \frac{|\{i : x_i^{(c)} = c\}|}{n_{\text{train}}}.
\label{eq:freq-enc}
\end{equation}

\textbf{Преобразование числовых признаков.} Лог-преобразование
тяжелохвостых признаков \(\tilde{x} = \log(1 + x)\) с последующим
robust-scaling:
\begin{equation}
\tilde{x}' = \frac{\tilde{x} - \operatorname{med}(\tilde{x})}
                  {\operatorname{IQR}(\tilde{x})}.
\label{eq:robust-scaler}
\end{equation}
Робастные оценки устойчивы к выбросам --- критично для NIDS,
где аномалии формируют ``хвосты'' распределения.

\textbf{Скользящие окна.} Принят размер \(W = 32\) с шагом
\(S = 8\) (75\% перекрытие). Метка окна формируется
по последнему flow (\emph{last}-агрегация) с учётом эмпирически
обоснованной структуры UNSW-NB15.

\subsubsection{Унификация признаков с CICIDS2017 для cross-evaluation}

Для проверки переносимости методики (гипотеза $H_0^{\text{transfer}}$)
вспомогательный датасет CICIDS2017~\cite{sharafaldin2018toward}
формируется утилитой CICFlowMeter и содержит более 80 flow-признаков.
Часть из них функционально эквивалентна полям UNSW-NB15.
Подмножество общих признаков, принятое в работе, представлено в
\tabref{tab:common-features}.

\begin{table}[H]
\centering
\caption{Подмножество признаков, общих для UNSW-NB15 и CICIDS2017,
принятое для cross-evaluation}
\label{tab:common-features}
\renewcommand{\arraystretch}{1.2}
\begin{tabularx}{\linewidth}{@{}lXX@{}}
\toprule
\textbf{Семантика} & \textbf{UNSW-NB15} & \textbf{CICIDS2017 (CICFlowMeter)} \\
\midrule
Длительность потока        & \emph{dur}      & Flow Duration \\
Протокол транспорта        & \emph{proto}    & Protocol \\
Прикладной сервис          & \emph{service}  & Destination Port (как proxy) \\
Байты от источника         & \emph{sbytes}   & Total Length of Fwd Packets \\
Байты к источнику          & \emph{dbytes}   & Total Length of Bwd Packets \\
Пакеты от источника        & \emph{spkts}    & Total Fwd Packets \\
Пакеты к источнику         & \emph{dpkts}    & Total Backward Packets \\
Скорость источника         & \emph{sload}    & Fwd Packets/s \\
Скорость получателя        & \emph{dload}    & Bwd Packets/s \\
Межпакетный интервал, src  & \emph{sinpkt}   & Fwd IAT Mean \\
Межпакетный интервал, dst  & \emph{dinpkt}   & Bwd IAT Mean \\
Джиттер, src               & \emph{sjit}     & Fwd IAT Std \\
Джиттер, dst               & \emph{djit}     & Bwd IAT Std \\
TCP-timing-сигнатуры       & \emph{synack, ackdat} & Flow IAT-агрегаты \\
\bottomrule
\end{tabularx}
\end{table}

Признаки, не имеющие прямого аналога (\emph{ct\_*}-агрегаты UNSW-NB15,
специфические признаки CICFlowMeter), исключаются из набора
cross-evaluation. В режиме \emph{train on UNSW, test on CICIDS2017}
вход модели сужается до 14 общих признаков; в режиме обучения только на
UNSW-NB15 используется полный набор из 60 признаков, формируемый
\texttt{Preprocessor}-ом.

\subsubsection{Управление дисбалансом классов}

Применяется \textbf{focal loss}~\cite{lin2017focal}:
\begin{equation}
\mathcal{L}_{\text{focal}} =
-\frac{1}{n}\sum_{i=1}^{n}
  \alpha_{y_i}\,(1-p_{y_i})^{\gamma}\,\log p_{y_i},
\label{eq:focal}
\end{equation}
с параметрами \(\alpha = 0{,}25\), \(\gamma = 2{,}0\). Для классических
baselines дополнительно применяется class re-weighting
\(w_c = n / (|\mathcal{C}|\,n_c)\).

\subsection{Архитектура ИИ-злоумышленника}

Реализуется гибридная архитектура \emph{templates + FSM + drift}:

\begin{enumerate}
    \item \textbf{Библиотека шаблонов} (семь параметризованных
    шаблонов: \texttt{ddos\_synflood, port\_scan, brute\_force,
    exploit\_payload, data\_exfil, internal\_recon, worm\_lateral})
    плюс шаблон \texttt{normal} для генерации фонового трафика.
    \item \textbf{Конечный автомат} с состояниями
    \[
    \mathcal{S} = \{\text{NORMAL}\} \cup
    \{\text{ATTACK}(c)\}_{c \in \mathcal{C}} \cup
    \{\text{DRIFT}(k)\}_{k \in \mathcal{K}}.
    \]
    Политика \(\pi\colon \mathcal{S} \to \Delta(\mathcal{S})\)
    задаётся вероятностями переходов и минимальными длительностями
    состояний.
    \item \textbf{Drift-инжектор} с тремя типами дрейфа:
    \emph{covariate, prior, concept}. Контроль интенсивности ---
    через индексы PSI и MMD относительно эталонного окна нормы.
\end{enumerate}

\subsection{Протокол экспериментов}

Сформирован протокол шести экспериментов, представленный в
\tabref{tab:experiments}.

\begin{table}[H]
\centering
\caption{Протокол экспериментов}
\label{tab:experiments}
\renewcommand{\arraystretch}{1.2}
\setlength{\tabcolsep}{4pt}
\small
\begin{tabularx}{\linewidth}{@{}lXXX@{}}
\toprule
\textbf{ID} & \textbf{Цель} & \textbf{Метрики} & \textbf{Критерий принятия} \\
\midrule
E1 &
Сравнение моделей-кандидатов на UNSW-NB15; одинаковый препроцессинг
и пороги &
Precision, Recall, F1, MCC, PR-AUC, ROC-AUC; $\sigma_{F1}$ &
Статистически значимое превосходство (\(p < 0{,}05\), Wilcoxon)
над ансамблем baselines \\
\addlinespace
E2 &
Анализ ошибок (FP/FN) лучшей модели; разрезы по \emph{attack\_cat} &
confusion matrix, FN/FP по классам, ROC по классам &
Идентификация классов с пониженной полнотой \\
\addlinespace
E3 &
Калибровка порога; температурное масштабирование &
ECE, FP-rate vs Recall, FP-per-time, TTD &
Рабочая точка по NFR-2 \\
\addlinespace
E4 &
Тестирование с ИИ-злоумышленником на стенде &
Precision, Recall, F1, TTD, FP-per-time; per-state recall &
$\Delta F_1 \leq 0{,}10$ при agent-NORMAL+ATTACK; TTD$\leq W$ \\
\addlinespace
E5 &
Drift-устойчивость; PSI/MMD-инжектирование &
PSI, MMD, $\Delta F_1$ vs PSI; время восстановления &
$\Delta F_1 \leq 0{,}10$ при PSI $\leq 0{,}25$;
срабатывание drift-monitor \\
\addlinespace
E6 &
Эксплуатационная производительность &
inference latency, throughput, CPU/RAM &
NFR-3, NFR-4 \\
\bottomrule
\end{tabularx}
\end{table}

Для каждого эксперимента применяются: множественные seed (\(n \geq 5\)
для лидирующих моделей), bootstrap CI~95\,\%, парный критерий
Уилкоксона при сравнении пар моделей, фиксированные YAML-конфиги.

\subsection*{Выводы по главе 2}

\begin{enumerate}
    \item Систематизированы опубликованные NIDS-результаты восьми
    типов архитектур с явной привязкой к первоисточникам.
    Установлены три методологических ограничения литературного
    сравнения, определяющие необходимость собственного эксперимента в
    равных условиях.
    \item Обоснован выбор CNN-LSTM как proposed-архитектуры по
    четырём аргументам: соответствие природе данных, опубликованные
    результаты, воспроизводимость, эксплуатационная стоимость.
    Альтернативы и контрольные модели зафиксированы.
    \item Описана теоретическая модель данных и pipeline обработки
    UNSW-NB15: схема признаков, очистка, encoding, log-преобразование,
    robust-scaling, скользящие окна (\(W=32\), \(S=8\), last-агрегация),
    focal loss для управления дисбалансом классов.
    \item Сформирована архитектура ИИ-злоумышленника:
    7 шаблонов + FSM + drift-инжектор; описаны три типа дрейфа.
    \item Сформирован протокол шести экспериментов E1--E6 с
    количественными критериями принятия и статистической методологией.
\end{enumerate}
